\chapter{周环反应}\label{ch:b3:pericyclic}

阳光照到皮肤上，使 7-脱氢胆固醇的一个环在几分之
一秒内打开；随后，体温使一个氢原子沿七个碳原子移动，
产物为维生素 D$_3$。两步均无中间体：化学键断裂
与形成同时进行，沿一圈原子排列，经由单一过渡态。
第一步只在光照下发生，第二步只在加热时发生，两步各
得到一种立体异构体。本章由所涉分子轨道的对称性
出发解释其原因：罗伯特·伍德沃德与罗尔德·
霍夫曼所发现的规则可预测这类反应能否发生及会得到什么。

\begin{recall}
本科二年级分册给出了多烯的休克尔系数、芳香性与
反芳香性环、HOMO、LUMO 与前线轨道、协同反应，以及
双烯体与亲双烯体的狄尔斯–阿尔德反应及其内型规则；本科一年级分册
给出 \textit{Z}/\textit{E} 描述符。在\Cref{ch:b3:point-groups}中定义了 $C_2$
轴和\omterm{def:b3:point-groups:mirror}{镜面}，在\cref{ch:b3:photochemistry}中讨论了激发态与
[2+2] 光\omterm{def:b3:pericyclic:families}{环加成}。
\end{recall}

\begin{omfigure}
\includegraphics[width=0.55\linewidth]{images/book4/ai/b3-26-terrace.jpg}

{\small 晴日露台上阅读：到达皮肤的紫外线开始
通过周环开环反应合成维生素 D$_3$。}
\end{omfigure}

\section{四类反应}

\begin{definition}[周环反应]\label{def:b3:pericyclic:pericyclic}
称\emph{周环反应}\index{周环反应}为一类协同反应，
其过渡态是相互作用轨道构成的环状排列，
所有化学键在环上同时形成与断裂。
\end{definition}

\begin{definition}[周环反应的类型]\label{def:b3:pericyclic:families}
在\emph{电环化反应}\index{电环化反应}中，共轭多烯
通过两端之间形成一个 $\sigma$ 键而闭合成环，或
其逆过程。在\emph{环加成}\index{环加成}中，两个 $\pi$ 体系
通过形成两个 $\sigma$ 键而连接成环；$[m + n]$ 环加成使
分别含 $m$ 个和 $n$ 个原子（或电子）的组分相连。在\emph{σ 迁移
重排}\index{σ 迁移重排}中，一个 $\sigma$ 键沿 $\pi$
体系迁移；在 $[i, j]$ 迁移中，键由两个片段的 1 号与 1 号原子
迁移到 $i$ 号与 $j$ 号原子上。\emph{螯键反应}\index{螯键反应}是
使同一原子上的两个键同时形成或断裂的反应，如二氧化硫
从磺烯中脱去。在\emph{烯反应}\index{烯反应}中，烯烃
带烯丙位氢，加成到一个 $\pi$ 键上，该氢原子随之迁移。
\end{definition}

\begin{definition}[同面与异面]\label{def:b3:pericyclic:faciality}
某组分以\emph{同面}\index{同面}方式反应，是指其两个新键形成
（或两个旧键断裂）时位于其 $\pi$ 体系的同一面（或
在 $\sigma$ 键处构型保持）；而以\emph{异面}\index{异面}方式反应时，
它们位于相反的两面（或构型翻转）。
\end{definition}

\section{电环化反应}

\begin{definition}[顺旋与对旋]\label{def:b3:pericyclic:rotation-modes}
在电环化闭环中，两个末端基团绕
各自与链相连的键转动。两者沿同一方向转动时，此旋转称为\emph{顺旋}\index{顺旋}
（沿链看二者均为顺时针），两者沿相反方向转动时则称为
\emph{对旋}\index{对旋}。
\end{definition}

\begin{theorem}[电环化反应]\label{thm:b3:pericyclic:electrocyclic}
含 $4n$ 个 $\pi$ 电子的多烯的热\omterm{def:b3:pericyclic:families}{电环化反应}
为\omterm{def:b3:pericyclic:rotation-modes}{顺旋}；含 $4n + 2$ 个 $\pi$ 电子者为\omterm{def:b3:pericyclic:rotation-modes}{对旋}。
\end{theorem}

\begin{proof}
新的 $\sigma$ 键由 HOMO 的末端 p 轨道形成，该轨道在
热反应中占据能量最高的电子；只有当相互转向的
瓣符号相同时，该组合才是成键的。由休克尔公式
（见\cref{ch:b3:solid-state}）可知，轨道 $k$ 在 $N$ 原子链中的系数
正比于 $\sin(jk\pi/(N + 1))$，于是两端 $j = 1$ 与 $j = N$ 的系数为
$\sin(k\pi/(N + 1))$ 与 $\sin(Nk\pi/(N + 1)) = (-1)^{k+1}\sin(k\pi/(N + 1))$：$k$ 为奇数时符号相同，
$k$ 为偶数时符号相反。含 $N$ 个原子和 $N$ 个电子时，HOMO 为
$k = N/2$。当 $N = 4n$ 时，$k = 2n$ 为偶数，两端的上部瓣
符号相反，两端须同向转动，各自将符号相配的瓣
转向内侧，为\omterm{def:b3:pericyclic:rotation-modes}{顺旋}。当 $N = 4n + 2$ 时，$k = 2n + 1$ 为奇数，上部瓣
符号相同，使它们相互靠近的转动为\omterm{def:b3:pericyclic:rotation-modes}{对旋}。
\end{proof}

\begin{omfigure}
\begin{tikzpicture}[font=\small, >=Stealth]
  % butadiene HOMO psi2: 0.60 0.37 -0.37 -0.60 (lobes scaled x1.6)
  \foreach \x/\c in {0/0.60, 1/0.37, 2/-0.37, 3/-0.60} {
    \pic at (\x,0) {omp={1.6*\c}{90}};
    \fill (\x,0) circle (1.2pt);
  }
  \draw[black!60] (0,0) -- (3,0);
  \draw[->, thick, omThm] (-0.45,0.55) arc[start angle=130, end angle=20, radius=0.7];
  \draw[->, thick, omThm] (2.55,0.55) arc[start angle=130, end angle=20, radius=0.7];
  \node[font=\scriptsize, align=center] at (1.5,-1.3) {丁二烯，$\psi_2$（HOMO，4 个电子）：\\ 两端符号相反：\textbf{顺旋}};
  \begin{scope}[xshift=6.2cm]
  % hexatriene HOMO psi3: 0.52 0.23 -0.42 -0.42 0.23 0.52
  \foreach \x/\c in {0/0.52, 0.8/0.23, 1.6/-0.42, 2.4/-0.42, 3.2/0.23, 4.0/0.52} {
    \pic at (\x,0) {omp={1.6*\c}{90}};
    \fill (\x,0) circle (1.2pt);
  }
  \draw[black!60] (0,0) -- (4.0,0);
  \draw[->, thick, omThm] (-0.45,0.55) arc[start angle=130, end angle=20, radius=0.7];
  \draw[->, thick, omThm] (4.45,0.55) arc[start angle=50, end angle=160, radius=0.7];
  \node[font=\scriptsize, align=center] at (2.0,-1.3) {己三烯，$\psi_3$（HOMO，6 个电子）：\\ 两端符号相同：\textbf{对旋}};
  \end{scope}
\end{tikzpicture}

{\small 丁二烯与己三烯的 HOMO，其瓣
按各自的休克尔系数缩放（蓝色与橙色代表两种符号），沿链绘出。
若两端之间要形成键，则相互转向的瓣
必须符号相同：在丁二烯中两端均顺时针转动，在己三烯中
两端则沿相反方向转动（红色箭头）。}
\end{omfigure}

该规则具有立体专一性。（\textit{E},\textit{E}）-己-2,4-二烯加热时，会\omterm{def:b3:pericyclic:rotation-modes}{顺旋}闭环
生成\textit{反式}-3,4-二甲基环丁烯；实际上，张力较大的
环丁烯才是开环的一方，而\textit{顺式}-3,4-二甲基环丁烯开环
只\omterm{def:b3:pericyclic:rotation-modes}{顺旋}生成（\textit{E},\textit{Z}）-己-2,4-二烯。（\textit{E},\textit{Z},\textit{E}）-辛-2,4,6-三烯
\omterm{def:b3:pericyclic:rotation-modes}{对旋}闭环，生成\textit{顺式}-5,6-二甲基环己-1,3-二烯。

\begin{proposition}[光化学电环化反应]\label{prop:b3:pericyclic:photo-electrocyclic}
在光照下，规则恰好相反：$4n$ 个电子，\omterm{def:b3:pericyclic:rotation-modes}{对旋}；$4n + 2$，\omterm{def:b3:pericyclic:rotation-modes}{顺旋}。
\end{proposition}

\begin{proof}
吸收光子后，电子由 HOMO（$k = N/2$）跃迁至 LUMO（$k = N/2 + 1$）
（见\cref{ch:b3:photochemistry}），此时它成为最高占据轨道，并
支配两端。其 $k$ 具有相反的奇偶性，因此其末端系数的相对符号
发生反转，旋转方向也随之改变。
\end{proof}

\begin{method}[预测电环化反应的立体化学]\label{met:b3:pericyclic:predict-stereo}
\begin{enumerate}
  \item 数出开链多烯的 $\pi$ 电子数（$4n$ 或 $4n + 2$）；注意是加热还是
        光照。
  \item 判断旋转模式：热反应中 $4n$ \omterm{def:b3:pericyclic:rotation-modes}{顺旋}，$4n + 2$ \omterm{def:b3:pericyclic:rotation-modes}{对旋}；光化学反应则相反。
  \item 画出两端及其取代基，取能闭环的构象（s-顺式）；
        按所选模式将两端旋转 90°。
  \item 判断两个取代基落在环的同一面
        （\textit{顺式}）还是相对的两面（\textit{反式}）。对于开环反应，
        则反向进行同样的分析。
\end{enumerate}
\end{method}

\section{轨道相关图}

伍德沃德与霍夫曼最初的论证考察所有轨道，而不只是
HOMO。沿\omterm{def:b3:pericyclic:rotation-modes}{顺旋}路径，分子保持一个 $C_2$ 轴；沿\omterm{def:b3:pericyclic:rotation-modes}{对旋}路径，
则保持\omterm{def:b3:point-groups:mirror}{镜面} $\sigma$。反应物与产物的
每个轨道相对于所保持的\omterm{def:b3:point-groups:operation}{对称元素}，
均为对称（S）或反对称（A）。

\begin{definition}[轨道相关图]\label{def:b3:pericyclic:orbital-correlation}
依据反应路径上保持的\omterm{def:b3:point-groups:operation}{对称元素}，将反应物的每个轨道与产物中
对称性相同的轨道相连，
连接时按能量顺序进行，即最低的 S 与最低的 S 相连，依此类推，
称为\emph{轨道相关图}\index{轨道相关图}。
\end{definition}

\begin{proposition}[轨道对称性守恒]\label{prop:b3:pericyclic:correlation}
若反应物中每个被占据的轨道都与产物中一个被占据的轨道相关联，
则热反应是允许的；若某个被占据的轨道与一个空轨道相关联，
则该反应被禁阻，且势垒很高。
\end{proposition}

\begin{proof}[论证]
沿反应路径，轨道保持其对称性，能量连续变化；
同对称性的轨道不能相交（它们会混合并相互排斥），
不同对称性的轨道则可以相交。当反应物中一个双占据轨道
必须变为产物中的反键轨道时，反应物的基态
将导致产物的双激发态，基态
能量急剧上升，直到反应路径后段，同对称性的态
发生混合，故势垒很高。
\end{proof}

\begin{omfigure}
\begin{tikzpicture}[font=\scriptsize]
  \foreach \px/\title/\el in {0/对旋：保持 $\sigma$/s, 7.6/顺旋：保持 $C_2$/c} {
    \begin{scope}[xshift=\px cm]
    \node[font=\small\bfseries] at (2.0,5.1) {\title};
    \node at (0,4.7) {丁二烯};
    \node at (4.0,4.7) {环丁烯};
    \end{scope}
  }
  % ---- disrotatory (mirror): butadiene psi1 S, psi2 A, psi3 S, psi4 A;
  %      cyclobutene sigma S, pi S, pi* A, sigma* A
  \pic (b1) at (0,0.6) {omlevel={ud}{1.0}}; \node[anchor=east] at (-0.6,0.6) {$\psi_1$ S};
  \pic (b2) at (0,1.8) {omlevel={ud}{1.0}}; \node[anchor=east] at (-0.6,1.8) {$\psi_2$ A};
  \pic (b3) at (0,3.0) {omlevel={}{1.0}}; \node[anchor=east] at (-0.6,3.0) {$\psi_3$ S};
  \pic (b4) at (0,4.1) {omlevel={}{1.0}}; \node[anchor=east] at (-0.6,4.1) {$\psi_4$ A};
  \pic (c1) at (4.0,0.1) {omlevel={ud}{1.0}}; \node[anchor=west] at (4.6,0.1) {$\sigma$ S};
  \pic (c2) at (4.0,1.3) {omlevel={ud}{1.0}}; \node[anchor=west] at (4.6,1.3) {$\pi$ S};
  \pic (c3) at (4.0,3.4) {omlevel={}{1.0}}; \node[anchor=west] at (4.6,3.4) {$\pi^*$ A};
  \pic (c4) at (4.0,4.4) {omlevel={}{1.0}}; \node[anchor=west] at (4.6,4.4) {$\sigma^*$ A};
  \draw[omcorr] (b1-r) -- (c1-l);
  \draw[omcorrforbidden] (b2-r) -- (c3-l);
  \draw[omcorr] (b3-r) -- (c2-l);
  \draw[omcorr] (b4-r) -- (c4-l);
  % ---- conrotatory (C2): butadiene psi1 A, psi2 S, psi3 A, psi4 S;
  %      cyclobutene sigma S, pi A, pi* S, sigma* A
  \begin{scope}[xshift=7.6cm]
  \pic (d1) at (0,0.6) {omlevel={ud}{1.0}}; \node[anchor=east] at (-0.6,0.6) {$\psi_1$ A};
  \pic (d2) at (0,1.8) {omlevel={ud}{1.0}}; \node[anchor=east] at (-0.6,1.8) {$\psi_2$ S};
  \pic (d3) at (0,3.0) {omlevel={}{1.0}}; \node[anchor=east] at (-0.6,3.0) {$\psi_3$ A};
  \pic (d4) at (0,4.1) {omlevel={}{1.0}}; \node[anchor=east] at (-0.6,4.1) {$\psi_4$ S};
  \pic (e1) at (4.0,0.1) {omlevel={ud}{1.0}}; \node[anchor=west] at (4.6,0.1) {$\sigma$ S};
  \pic (e2) at (4.0,1.3) {omlevel={ud}{1.0}}; \node[anchor=west] at (4.6,1.3) {$\pi$ A};
  \pic (e3) at (4.0,3.4) {omlevel={}{1.0}}; \node[anchor=west] at (4.6,3.4) {$\pi^*$ S};
  \pic (e4) at (4.0,4.4) {omlevel={}{1.0}}; \node[anchor=west] at (4.6,4.4) {$\sigma^*$ A};
  \draw[omcorr] (d1-r) -- (e2-l);
  \draw[omcorr] (d2-r) -- (e1-l);
  \draw[omcorr] (d3-r) -- (e4-l);
  \draw[omcorr] (d4-r) -- (e3-l);
  \end{scope}
\end{tikzpicture}

{\small 丁二烯 $\rightleftharpoons$ 环丁烯的\omterm{def:b3:pericyclic:orbital-correlation}{轨道相关图}。左图，
\omterm{def:b3:pericyclic:rotation-modes}{对旋}路径保持\omterm{def:b3:point-groups:mirror}{镜面}：被占据的 $\psi_2$（A）与
环丁烯的空 $\pi^*$（A）相关（红色虚线）：热反应禁阻。右图，
\omterm{def:b3:pericyclic:rotation-modes}{顺旋}路径保持 $C_2$ 轴：被占据轨道与被占据轨道相连：
热反应允许。}
\end{omfigure}

\begin{method}[绘制相关图]\label{met:b3:pericyclic:correlation-diagram}
\begin{enumerate}
  \item 找出平分形成键与断裂键的\omterm{def:b3:point-groups:operation}{对称元素}，
        并要求它们在所选路径上始终保持。
  \item 列出参与反应的反应物与产物轨道，按能量排序
        并对每个\omterm{def:b3:point-groups:operation}{对称元素}标注 S 或 A。
  \item 将一侧每种对称类型的最低轨道，与另一侧同类型的最低轨道相连，
        再依次连接下一组，
        同一对称性的两条连线绝不交叉。
  \item 检查被占据轨道：若全部连到被占据轨道，则允许；若有一个连到空轨道，
        则禁阻。
\end{enumerate}
\end{method}

\begin{definition}[状态相关图]\label{def:b3:pericyclic:state-correlation}
\emph{状态相关图}\index{状态相关图}沿反应路径，将反应物与产物的
电子态相互连接。每个电子态由轨道组态构成，
并以其总对称性加以标记；
相同对称性的电子态互不交叉。
\end{definition}

对于禁阻的\omterm{def:b3:pericyclic:rotation-modes}{对旋}路径，丁二烯的基态（$\psi_1^2\psi_2^2$）
与环丁烯的双激发态（$\sigma^2\pi^{*2}$）相关联，而第一激发态
（$\psi_1^2\psi_2\psi_3$）则与环丁烯的第一激发态（$\sigma^2\pi\pi^*$）相关联：从激发态出发，
\omterm{def:b3:pericyclic:rotation-modes}{对旋}路径是下坡的。这正是
\cref{prop:b3:pericyclic:photo-electrocyclic}所述光化学逆反应在电子态上的体现。

\section{环加成}

\begin{theorem}[环加成]\label{thm:b3:pericyclic:cycloaddition}
两组分均以\omterm{def:b3:pericyclic:faciality}{同面}方式反应的热 $[m + n]$ \omterm{def:b3:pericyclic:families}{环加成}，
当 $m + n = 4q + 2$ 个电子参与时为允许，$m + n = 4q$ 时为禁阻；光照下
规则反转。
\end{theorem}

\begin{proof}
两个新键都形成于一个组分的 HOMO 与另一组分的 LUMO 之间，
二者以端部瓣相对。两组分均\omterm{def:b3:pericyclic:faciality}{同面}反应时，
两个键都成键，仅当 HOMO 与 LUMO 的端部系数
的相对符号相同。对于含 $N$ 个电子、位于 $N$ 个原子上的组分，
其 HOMO 的 $k = N/2$，LUMO 的 $k = N/2 + 1$；端部系数
在奇数 $k$ 时符号相同，偶数 $k$ 时符号相反。在 [4+2] 情形中，二烯的 HOMO（$k = 2$）
与乙烯的 LUMO（$k = 2$）端部符号均相反，二者相匹配。在
[2+2] 情形中，一个乙烯的 HOMO（$k = 1$，符号相同）与另一个
乙烯的 LUMO（$k = 2$，符号相反）相遇：一端成键，一端反键。一般地，$m$ 组分的 HOMO
与 $n$ 组分的 LUMO 相匹配，当 $m/2$ 与 $n/2 + 1$ 的
奇偶性相同时，即 $(m + n)/2$ 为奇数，$m + n = 4q + 2$。在光照下，激发组分的单占据 LUMO
起着其 HOMO 的作用，奇偶性条件
反转。
\end{proof}

\begin{omfigure}
\begin{tikzpicture}[font=\scriptsize]
  % [4+2]: diene HOMO above, ethene LUMO (pi*) below, facing lobes match at both ends
  \foreach \x/\c in {0/0.60, 1/0.37, 2/-0.37, 3/-0.60} {
    \pic at (\x,1.6) {omp={1.4*\c}{90}}; \fill (\x,1.6) circle (1.1pt);
  }
  \draw[black!60] (0,1.6) -- (3,1.6);
  \foreach \x/\c in {0/-0.71, 3/0.71} {
    \pic at (\x,-0.4) {omp={1.4*\c}{90}}; \fill (\x,-0.4) circle (1.1pt);
  }
  \draw[black!60] (0,-0.4) -- (3,-0.4);
  \draw[omMeth, densely dashed, thick] (-0.25,0.75) -- (-0.25,0.35) (3.25,0.75) -- (3.25,0.35);
  \node[omMeth] at (-0.25,0.55) [left] {成键};
  \node[omMeth] at (3.25,0.55) [right] {成键};
  \node[font=\small] at (1.5,2.95) {[4+2]：二烯 HOMO};
  \node[font=\small] at (1.5,-1.75) {乙烯 LUMO：允许};
  \begin{scope}[xshift=7.2cm]
  \foreach \x/\c in {0/0.71, 3/0.71} {
    \pic at (\x,1.6) {omp={1.4*\c}{90}}; \fill (\x,1.6) circle (1.1pt);
  }
  \draw[black!60] (0,1.6) -- (3,1.6);
  \foreach \x/\c in {0/-0.71, 3/0.71} {
    \pic at (\x,-0.4) {omp={1.4*\c}{90}}; \fill (\x,-0.4) circle (1.1pt);
  }
  \draw[black!60] (0,-0.4) -- (3,-0.4);
  \draw[omMeth, densely dashed, thick] (-0.25,0.75) -- (-0.25,0.35);
  \draw[omThm, densely dashed, thick] (3.25,0.75) -- (3.25,0.35);
  \node[omMeth] at (-0.25,0.55) [left] {成键};
  \node[omThm] at (3.25,0.55) [right] {反键};
  \node[font=\small] at (1.5,2.95) {[2+2]：乙烯 HOMO};
  \node[font=\small] at (1.5,-1.75) {乙烯 LUMO：禁阻};
  \end{scope}
\end{tikzpicture}

{\small \omterm{def:b3:pericyclic:faciality}{同面}\omterm{def:b3:pericyclic:families}{环加成}中相互面对的前线轨道（上方组分的下瓣与下方组分的上瓣相对，
瓣的大小取自休克尔系数）。
在 [4+2] 反应中，两个新键都成键；
在 [2+2] 反应中，一端成键，另一端反键。}
\end{omfigure}

热反应中，两个烯烃不能经一步协同反应生成环丁烷；在
光照下则可以（见\cref{ch:b3:photochemistry}）。烯酮是一个例外，恰恰
证明了规则：其 C=O 的 $\pi^*$ 轨道与 C=C 的轨道垂直，使其
与烯烃以 [$\pi2_s$ + $\pi2_a$] 方式反应，一个组分\omterm{def:b3:pericyclic:faciality}{同面}参与，
烯酮\omterm{def:b3:pericyclic:faciality}{异面}参与，热反应下生成环丁酮。

\begin{definition}[1,3-偶极环加成]\label{def:b3:pericyclic:dipolar}
\emph{1,3-偶极子}\index{1,3-偶极子}是含四个电子的三原子 $\pi$ 体系，
其形式电荷标在两端或中间，例如叠氮化物
$\ce{R-N=N+=N-}$、腈氧化物或臭氧分子。
\emph{1,3-偶极环加成}\index{1,3-偶极环加成}是其与烯烃或炔烃发生的
[3+2] \omterm{def:b3:pericyclic:families}{环加成}，属六电子反应，生成五元环。
\end{definition}

叠氮化物与端炔的铜催化反应，生成一种
1,2,3-三唑的区域异构体，是点击化学中可靠的连接反应。
臭氧解始于臭氧与烯烃的
\omterm{def:b3:pericyclic:dipolar}{1,3-偶极环加成}。

\section{σ 迁移重排与一般规则}

\begin{proposition}[氢的迁移]\label{prop:b3:pericyclic:sigmatropic-h}
氢沿多烯的热 $[1, j]$ 迁移，当
过渡态含 $4q + 2$ 个电子（$[1,5]$）时允许\omterm{def:b3:pericyclic:faciality}{同面}进行，
含 $4q$ 个电子（$[1,3]$、$[1,7]$）时必须\omterm{def:b3:pericyclic:faciality}{异面}进行。
\end{proposition}

\begin{proof}
将过渡态视为氢原子（1s）在一个由 $j$ 个原子、含 $j$ 个电子的
自由基；氢架在 SOMO 的两个末端之间，
$k = (j + 1)/2$。当 $j = 5$ 时，$k = 3$ 为奇数，两末端的瓣在同一面上符号相同，
1s 轨道可在该面同时与两端成键：\omterm{def:b3:pericyclic:faciality}{同面}迁移。对于 $j = 3$ 与 $j = 7$，
$k = 2$ 与 $4$ 为偶数，同号瓣分处两端的相反两面，
氢必须从一面转到另一面：\omterm{def:b3:pericyclic:faciality}{异面}迁移。$[1,3]$ 迁移
因链太短而无法跨越；$[1,7]$ 迁移则可沿螺旋形多烯进行。
\end{proof}

\begin{definition}[科普重排与克莱森重排]\label{def:b3:pericyclic:cope-claisen}
\emph{科普重排}\index{科普重排}是己-1,5-二烯的 [3,3] σ 迁移重排；
\emph{克莱森
重排}\index{克莱森重排}是烯丙基乙烯基醚或
烯丙基芳基醚的重排，生成 $\gamma,\delta$-不饱和\omterm{def:b3:homogeneous-catalysis:carbonylation}{羰基化}合物或邻位-烯丙基苯酚。
\end{definition}

\begin{omfigure}
\begin{tikzpicture}[font=\scriptsize, scale=1.1]
  \foreach \px/\three/\title in {0/C/科普, 5.0/O/克莱森} {
    \begin{scope}[xshift=\px cm]
      % chair: atoms at 60-degree steps, alternately 0.25 above and below the
      % mean plane, seen from slightly above (y' = 0.35 y + z), scaled
      \begin{scope}[shift={(1.3,0.2)}, scale=1.15]
      \coordinate (a1) at (-1.25,-0.25); \coordinate (a2) at (-0.625,-0.129); \coordinate (a3) at (0.625,-0.629);
      \coordinate (a4) at (1.25,0.25); \coordinate (a5) at (0.625,0.129); \coordinate (a6) at (-0.625,0.629);
      \end{scope}
      \draw[thick] (a1) -- (a2) -- (a3);
      \draw[thick] (a4) -- (a5) -- (a6);
      \draw[omThm, thick, densely dashed] (a3) -- (a4);
      \draw[omMeth, thick, densely dashed] (a6) -- (a1);
      \foreach \i/\l in {1/1, 2/2, 4/4, 5/5, 6/6} { \fill (a\i) circle (1.3pt); }
      \node[fill=white, inner sep=0.6pt] at (a3) {\three};
      \node[anchor=east] at (a1) {1};
      \node[anchor=north] at (a2) {2};
      \node[anchor=north west] at (a3) {3};
      \node[anchor=west] at (a4) {4};
      \node[anchor=south] at (a5) {5};
      \node[anchor=south] at (a6) {6};
      \node[font=\small\bfseries] at (1.3,1.5) {\title};
    \end{scope}
  }
\end{tikzpicture}

{\small [3,3] σ 迁移重排的椅式过渡态：两个烯丙基片段彼此相对；
3--4 $\sigma$ 键断裂（红色），同时 1--6 键形成（绿色）。
在\omterm{def:b3:pericyclic:cope-claisen}{克莱森重排}中，原子 3 为醚氧原子，
产物为\omterm{def:b3:homogeneous-catalysis:carbonylation}{羰基化}合物。椅式构象决定了
新生成的手性碳原子的相对构型。}
\end{omfigure}

[3,3] 迁移是六电子、全部\omterm{def:b3:pericyclic:faciality}{同面}的过程，热允许；
它经由类似环己烷的椅式过渡态进行，
取代基倾向于占据其平伏键位置：（\textit{E},\textit{E}）与（\textit{Z},\textit{Z}）
双烯体分别给出相反的相对构型，可以预测。

\begin{theorem}[伍德沃德–霍夫曼规则]\label{thm:b3:pericyclic:woodward-hoffmann}
热\omterm{def:b3:pericyclic:pericyclic}{周环反应}被允许的条件是：$(4q + 2)_s$ 与
$(4r)_a$ 组分的总数为奇数；光化学\omterm{def:b3:pericyclic:pericyclic}{周环反应}则在该数为偶数时允许。
\end{theorem}

\begin{proof}[部分证明]
对上文讨论的每一族，计数与前线轨道分析一致：
热 [4+2] \omterm{def:b3:pericyclic:families}{环加成}含组分 $\pi4_s$ 与 $\pi2_s$，其中有一个 $(4q + 2)_s$
组分，数目为奇数，允许；[2+2] \omterm{def:b3:pericyclic:families}{环加成}含 $\pi2_s + \pi2_s$，共两个，为偶数，禁阻；而
$\pi2_s + \pi2_a$ 含一个，允许；环丁烯的\omterm{def:b3:pericyclic:rotation-modes}{顺旋}开环为 $\sigma2_s + \pi2_a$，
计一个，允许；\omterm{def:b3:pericyclic:faciality}{同面} $[1,5]$-H 迁移为 $\sigma2_s + \pi4_s$，其中只计 $\sigma2_s$：
共一个，允许。一般性证明适用于任意组分排列，
需在更高阶的课程中讲授。
\end{proof}

\begin{definition}[芳香过渡态与默比乌斯过渡态]\label{def:b3:pericyclic:mobius}
称\emph{芳香过渡态}\index{芳香过渡态}为周环过渡态中的
一组环状轨道，它与休克尔环一样，
在容纳 $4q + 2$ 个电子时趋于稳定。一个\emph{默比乌斯
过渡态}\index{默比乌斯过渡态}中，环上相邻轨道间的符号改变次数为奇数，
即含一次\omterm{def:b3:pericyclic:faciality}{异面}扭转；它
在容纳 $4q$ 个电子时趋于稳定。
\end{definition}

\begin{proposition}[芳香过渡态]\label{prop:b3:pericyclic:aromatic-ts}
热\omterm{def:b3:pericyclic:pericyclic}{周环反应}在其过渡态为\omterm{def:b3:pericyclic:mobius}{芳香过渡态}时被允许：
即休克尔拓扑且电子数为 $4q + 2$，或默比乌斯拓扑且电子数为 $4q$。
\end{proposition}

\begin{proof}[论证]
过渡态中的环状轨道组，其行为类似于
环状共轭分子。休克尔环的最低能级非简并，其上能级成对，
并以 $4q + 2$ 个电子填充（本科二年级分册）；含一次符号反转的环
其能级全部成对，以 $4q$ 个电子填充。填满的壳层意味着
稳定，即得到\omterm{def:b3:pericyclic:mobius}{芳香过渡态}与较低的势垒。符号改变的次数
即\omterm{def:b3:pericyclic:faciality}{异面}组分的个数，因此本规则与
伍德沃德–霍夫曼规则等价。
\end{proof}

\begin{method}[画图与计数组分]\label{met:b3:pericyclic:components}
\begin{enumerate}
  \item 标出形成和断裂的键；每个发生变化的 $\pi$ 体系或 $\sigma$ 键即为一个
        组分，以其类型和电子数标记（$\pi4$、$\pi2$、$\sigma2$、$\omega0$ 表示
        空轨道）。
  \item 依据几何构型，为每个组分选取 s 或 a（\omterm{def:b3:pericyclic:faciality}{同面}，或
        \omterm{def:b3:pericyclic:faciality}{异面}）。
  \item 计数 $(4q + 2)_s$ 与 $(4r)_a$ 组分：数目为奇数则热允许；数目为偶数则
        光化学允许。
  \item 以环上电子数交叉检验：休克尔型为 $4q + 2$，
        默比乌斯型为 $4q$。
\end{enumerate}
\end{method}

\begin{omfigure}
\begin{tikzpicture}[font=\scriptsize, >=Stealth]
  \foreach \px/\mode in {0/1, 4.6/2, 9.2/3} {
    \begin{scope}[xshift=\px cm]
      \foreach \n/\a in {10/210, 5/150, 6/90, 7/30, 8/-30, 9/-90} { \coordinate (v\n) at (\a:0.85); }
      \ifnum\mode=1
        \draw[thick] (v10) -- (v5) (v6) -- (v7) (v8) -- (v9);
        \draw[thick, double, double distance=1.6pt] (v5) -- (v6) (v7) -- (v8);
        \draw[thick, omThm] (v9) -- (v10);
        \draw[thick] (v10) -- ++(240:0.6) node[anchor=north] {$\ce{CH3}$};
        \node[font=\small\bfseries] at (0,1.55) {7-脱氢胆固醇};
      \fi
      \ifnum\mode=2
        \draw[thick, double, double distance=1.6pt] (v10) -- (v5) (v6) -- (v7) (v8) -- (v9);
        \draw[thick] (v5) -- (v6) (v7) -- (v8);
        \draw[thick] (v10) -- ++(240:0.6) node[anchor=north] (me) {$\ce{CH3}$};
        \draw[->, omThm, thick] ($(v10)+(240:0.75)+(0.35,0)$) to[bend right=40] ($(v9)+(-0.12,-0.12)$);
        \node[font=\small\bfseries] at (0,1.55) {前维生素 D$_3$};
      \fi
      \ifnum\mode=3
        \draw[thick] (v10) -- (v5) (v6) -- (v7) (v8) -- (v9);
        \draw[thick, double, double distance=1.6pt] (v5) -- (v6) (v7) -- (v8);
        \draw[thick, double, double distance=1.6pt] (v10) -- ++(240:0.6) node[anchor=north] {$\ce{CH2}$};
        \node[anchor=north west, inner sep=1pt] at (v9) {$\ce{CH2}$};
        \node[font=\small\bfseries] at (0,1.55) {维生素 D$_3$};
      \fi
      \foreach \n/\a in {10/210, 5/150, 6/90, 7/30, 8/-30, 9/-110} {
        \node[font=\tiny, black!60] at (\a:1.12) {\n};
      }
    \end{scope}
  }
  \draw[->, thick] (1.45,0) -- (3.0,0) node[midway, above, align=center] {$h\nu$\\ 6 e$^-$，顺旋};
  \draw[->, thick] (6.05,0) -- (7.6,0) node[midway, above, align=center] {加热，$[1,7]$-H\\ 异面};
\end{tikzpicture}

{\small 维生素 D 合成中的两个\omterm{def:b3:pericyclic:pericyclic}{周环反应}步骤，画在参与反应的原子上
（甾体编号用灰色标出；分子其余部分，即 A、C、D 环，
已略去）。光照使 B 环断开，断裂 9--10 键（红色），这是光化学六电子\omterm{def:b3:pericyclic:families}{电环化反应}，
\omterm{def:b3:pericyclic:rotation-modes}{顺旋}。在前维生素 D$_3$ 中，C19 甲基上的一个氢随后
迁移到 C9（红色箭头），穿过七原子螺旋三烯，
以\omterm{def:b3:pericyclic:faciality}{异面}方式进行。}
\end{omfigure}

\begin{inthelab}[克莱森重排]
将烯丙基芳基醚在氮气保护下，于高沸点
溶剂（如 1,2-二氯苯）中加热回流，或不用溶剂，在约 \qty{200}{\celsius} 下进行，并
用薄层色谱法跟踪反应进程，生成的\textit{邻}-烯丙基苯酚
比醚的极性更大。苯酚随后被萃入氢氧化钠水溶液中，
与中性杂质分离，再经酸化回收。
\end{inthelab}

\begin{safety}{\ghs[0.62cm]{GHS02}\ghs[0.62cm]{GHS04}\ghs[0.62cm]{GHS08}}
% ledger: ghs4:butadiene
1,3-丁二烯，一种常见的二烯：加压下为极易燃气体，可能导致
遗传缺陷和癌症。使用时仅限于封闭系统；教学中则
改用固体环丁砜烯，加热时经\omterm{def:b3:pericyclic:families}{螯键反应}释放丁二烯，
在封闭装置中代替它使用。
\end{safety}

\begin{history}[轨道对称性]
1965 年，罗伯特·伍德沃德在合成维生素 B$_{12}$ 时，发现了电环化步骤中令人费解的立体化学，
与罗尔德·霍夫曼共同发表了
轨道对称性守恒规则。福井谦一早在 1952 年便提出了
前线轨道；二人共同分享了 1981 年的
诺贝尔化学奖。伍德沃德曾因合成工作获 1965 年奖项，
于 1979 年去世。
\end{history}

\section{练习}

\begin{exercise}[$\star$]\label{exo:b3:pericyclic:1}
对下列过程进行分类：第尔斯–阿尔德反应；环丁烯的开环；\omterm{def:b3:pericyclic:cope-claisen}{科普
重排}；环丁砜烯失去 $\ce{SO2}$；臭氧对烯烃的加成；
以及环戊二烯中的 $[1,5]$-H 迁移。
\end{exercise}

\begin{exercise}[$\star$]\label{exo:b3:pericyclic:2}
预测闭环方式（\omterm{def:b3:pericyclic:rotation-modes}{顺旋}或\omterm{def:b3:pericyclic:rotation-modes}{对旋}），包括热反应与光化学反应中
己-1,3,5-三烯和丁-1,3-二烯的闭环。
\end{exercise}

\begin{exercise}[$\star$]\label{exo:b3:pericyclic:3}
\textit{反式}-3,4-二甲基环丁烯加热后生成何种产物？
\end{exercise}

\begin{exercise}[$\star$]\label{exo:b3:pericyclic:4}
为何乙烯加热时不能二聚生成环丁烷，而在光照下却能够二聚
（借助敏化剂或直接激发）？
\end{exercise}

\begin{exercise}[$\star\star$]\label{exo:b3:pericyclic:5}
画出己三烯 $\rightleftharpoons$ 环己-1,3-二烯在\omterm{def:b3:pericyclic:rotation-modes}{对旋}
途径下的\omterm{def:b3:pericyclic:orbital-correlation}{轨道相关图}，并由此作出结论。
\end{exercise}

\begin{exercise}[$\star\star$]\label{exo:b3:pericyclic:6}
在光照下，(\textit{E},\textit{Z},\textit{E})-辛-2,4,6-三烯闭环生成哪种二甲基环己二烯？
\end{exercise}

\begin{exercise}[$\star\star$]\label{exo:b3:pericyclic:7}
5-甲基环戊-1,3-二烯在室温下重排为 1-甲基
与 2-甲基环戊二烯。试用 $[1,5]$-H 迁移加以解释，并说明它为何是\omterm{def:b3:pericyclic:faciality}{同面}的。
\end{exercise}

\begin{exercise}[$\star\star$]\label{exo:b3:pericyclic:8}
数出烯酮与烯烃\omterm{def:b3:pericyclic:families}{环加成}中的组分数，并证明
该过程是热允许的。
\end{exercise}

\begin{exercise}[$\star\star$]\label{exo:b3:pericyclic:9}
烯丙基乙烯基醚 $\ce{CH2=CH-O-CH2-CH=CH2}$ 加热时
发生\omterm{def:b3:pericyclic:cope-claisen}{克莱森重排}。画出其椅式过渡态，并给出产物。
\end{exercise}

\begin{exercise}[$\star\star\star$]\label{exo:b3:pericyclic:10}
过渡态中八个电子呈环状排列，其中含一个\omterm{def:b3:pericyclic:faciality}{异面}
组分。该过渡态是否具有芳香性？该反应是否为热允许的？
\end{exercise}

\begin{exercise}[$\star\star\star$]\label{exo:b3:pericyclic:11}
叠氮化物与端炔的 \omterm{def:b3:pericyclic:dipolar}{1,3-偶极环加成}中，叠氮化物的最大
HOMO 系数位于末端氮原子 N3 上，端炔的最大
LUMO 系数位于其末端碳原子上（题目数据）。
预测无催化剂时占优的区域异构体，方法是将最大
系数相互配对，并加以讨论。
\end{exercise}

\begin{exercise}[$\star\star\star$]\label{exo:b3:pericyclic:12}
试用休克尔公式证明：$[1,7]$-H 迁移必为\omterm{def:b3:pericyclic:faciality}{异面}的，
而 $[1,5]$-H 迁移必为\omterm{def:b3:pericyclic:faciality}{同面}的。
\end{exercise}

% keep the heading with its problem box, which fills a page by itself
\clearpage
\enlargethispage{3\baselineskip}
\section{问题：阳光、皮肤与维生素 D}

\begin{problem}[{周末问题——阳光、皮肤与维生素 D：
7-脱氢胆固醇的光化学开环、热致氢迁移、
副产物，以及人体完成整个过程所需的时间}]
\label{pb:b3:pericyclic:1}
题目数据：在 \qty{37}{\celsius} 下，前维生素 D$_3$ 转化为维生素 D$_3$ 的
转化呈一级，$k = \qty{2.3e-2}{h^{-1}}$（忽略逆向转化）；其活化能为
\qty{85}{kJ/mol}.

\medskip\noindent\textbf{第一部分——开环过程。}
\begin{enumerate}
  \item 7-脱氢胆固醇中哪一根键断裂？形成了什么 $\pi$ 体系？
  \item 有多少个电子参与反应？该反应属于哪一类？
  \item 为什么需要光照？
  \item 光化学开环是\omterm{def:b3:pericyclic:rotation-modes}{顺旋}还是\omterm{def:b3:pericyclic:rotation-modes}{对旋}？
        请借助轨道加以论证。
  \item 为何前维生素 D$_3$ 的新中央双键必须是 \textit{Z} 构型？
  \item 为什么在体温下、黑暗中，同一开环过程不能热致发生？
\end{enumerate}

\medskip\noindent\textbf{第二部分——氢迁移。}
\begin{enumerate}[resume]
  \item 指出氢原子在哪两个原子之间迁移，并说明迁移类型。
  \item 过渡态中有多少个电子参与？
  \item 数一数各组分，判断该热反应
        是\omterm{def:b3:pericyclic:faciality}{同面}允许还是\omterm{def:b3:pericyclic:faciality}{异面}允许。
  \item 为何这里的\omterm{def:b3:pericyclic:faciality}{异面}迁移在几何上可行，
        而 $[1,3]$ 迁移不可行？
  \item 该过渡态是休克尔型还是默比乌斯型？是否具有芳香性？
  \item 为何该迁移不需要光照？
\end{enumerate}

\medskip\noindent\textbf{第三部分——副产物。}
\begin{enumerate}[resume]
  \item 前维生素 D$_3$ 在光照下可再次闭环，其方式是什么，
        以及它为何能生成 7-脱氢胆固醇的环式立体异构体（光甾醇）？
  \item 光照也可使中央的 \textit{Z} 双键转变为 \textit{E} 双键（速甾醇）。为何
        速甾醇不能发生 $[1,7]$-H 迁移？
  \item 为何长时间日照并不会产生越来越多的维生素 D？
  \item 为何开环很快，而氢迁移却很慢？
\end{enumerate}

\medskip\noindent\textbf{第四部分——体温下的动力学。}
\begin{enumerate}[resume]
  \item 计算前维生素 D$_3$ 在 \qty{37}{\celsius} 下的半衰期。
  \item 计算一天内的转化分数。
  \item 计算转化 90\,\% 所需的时间。
  \item 计算 \qty{20}{\celsius} 时的 $k$。
  \item 计算 \qty{20}{\celsius} 下转化 90\,\% 所需的时间。
  \item 皮肤温暖有何好处？
  \item 为何维生素 D 补充剂要通过照射并加热甾醇来制备，
        而不采用经典合成？
  \item 给出结论：前维生素 D$_3$ 在体温下转化 90\,\% 所需的时间。
\end{enumerate}
\end{problem}
